\section{Choosing the smaller base and completing the induction}
\label{sec:induction}

We now prove Theorem~\ref{thm:finite-volume}. The transfer proposition
requires a rational section without horizontal poles and adjoint rank
one in every nonnegative degree. The task of this section is to obtain
both conditions from failure of nonvanishing. A maximal linear system
will give a strictly smaller base; the same maximality will control all
ratios on its generic fiber.

\subsection{One fixed determinant with unbounded twists}

We first explain the passage from cohomology to a single line bundle.
The determinant construction is a standard way to extract effective
divisors from twisted differentials; compare the determinant method of
Lazi\'c--Peternell \cite[Lemma~4.1]{LP2018}. We include the
argument because its uniformity in arbitrarily far tails is needed
here, including when the numerical class of $L$ is zero.

\begin{lemma}\label{lem:determinant}
Let $Z$ be smooth connected projective of dimension $d>0$, with
$H^0(Z,\Omega_Z^q)=0$ for $1\le q\le d$. Let $L$ have a
semipositive metric with locally bounded weights, and let $\alpha$ be
a nonzero movable class satisfying $\mu_{\min,\alpha}(T_Z)\ge0$.
There are a line bundle $M$ with $M\cdot\alpha\le0$, strictly
increasing positive integers $m_i\to\infty$, and nonzero sections
\[
                     s_i\in H^0(Z,M+m_iL).
\]
The line $M$ is the determinant of a saturated subsheaf of
$\Omega_Z^p$ for some $0\le p\le d$.
\end{lemma}

\begin{proof}
Hodge symmetry gives $H^q(Z,\OO_Z)=0$ for $q>0$, and hence
$\chi(\OO_Z)=1$. Riemann--Roch makes $\chi(K_Z+mL)$ a polynomial
in the integer $m$, whose constant term is
\[
                  \chi(K_Z)=(-1)^d\chi(\OO_Z)=(-1)^d
\]
by Serre duality. It is therefore nonzero for all but finitely many
positive integers. Some fixed $q$ has $H^q(Z,K_Z+mL)\ne0$ for
unbounded positive $m$.

For each such $m$, the metric $h_L^m$ has trivial multiplier ideal:
its weight is locally bounded, so every holomorphic germ is locally
square integrable. Hard Lefschetz with multiplier ideals
\cite[Theorem~2.1.1]{DPS2001} gives a surjection
\[
 H^0(Z,\Omega_Z^{d-q}\otimes mL)
       \longrightarrow H^q(Z,K_Z+mL)
\]
by wedging with the $q$th power of a K\"ahler form. Choose nonzero
maps $-mL\to U$, where $U=\Omega_Z^{d-q}$, for an increasing
sequence of these exponents.

At the generic point of $Z$, take the subspace of $U_{\C(Z)}$
spanned by all chosen images after a given cutoff. These subspaces
form a decreasing chain in a finite-dimensional vector space. Their
dimensions eventually stabilize, so the subspaces themselves stabilize
to a nonzero subspace $W$. Let $\mathcal G\subset U$ be its saturated
extension. It is coherent: saturate the images of finitely many maps
spanning $W$. Every map in a sufficiently far tail factors through
$\mathcal G$, because $U/\mathcal G$ is torsion-free and the
composite vanishes generically.

Put $r=\operatorname{rank}\mathcal G$ and
$M=(\bigwedge^r\mathcal G)^{**}$. This is a line bundle on smooth
$Z$. For every cutoff, choose $r$ generically independent maps from
the tail. Their exterior product is a nonzero section of $M+mL$,
where $m$ is the sum of their exponents. As the cutoff tends to
infinity, so does $m$. The maps and their determinant extend across
codimension two into the line $M$. Selecting an increasing subsequence
gives the sections in the statement.

For $d-q>0$, antisymmetrization embeds $U$ in
$(\Omega_Z^1)^{\otimes(d-q)}$. The tensor slope formula and the
assumed tangent inequality give $\mu_{\max,\alpha}(U)\le0$,
and therefore $M\cdot\alpha\le0$. If $d-q=0$, then
$U=\OO_Z$ and its nonzero saturated subsheaf is $\OO_Z$ itself;
thus $M=\OO_Z$ and the same conclusion holds.
\end{proof}

\subsection{A maximal system controls every section ratio}

The next lemma is algebraic. It isolates the dimension drop and the
rank calculation from the analytic transfer. The nonpositive degrees
are used only to prevent the base from having the full dimension.

\begin{lemma}[Maximal base and adjoint rank]\label{lem:maximal-base}
Let $Z$ be smooth connected projective of positive dimension $d$.
Let $L,M$ be line bundles and $D\ge0$ an integral divisor. Suppose
there is a nonzero movable class $\alpha$ with
\[
 L\cdot\alpha=D\cdot\alpha=0,\qquad M\cdot\alpha\le0,
\]
and nonzero sections $s_i\in H^0(Z,M+m_iL)$ for strictly increasing
positive integers $m_i\to\infty$.
There is a smooth projective base $S$ with $\dim S<d$, together with
the prepared diagram of Lemma~\ref{lem:prepared-model} and an integer
$k>0$, for which:
\begin{enumerate}[label=\textup{(\roman*)}]
\item a nonzero rational section $\rho$ of $k\bar\pi^*L$ has no
      horizontal poles on $\bar Y$;
\item putting $J=-K_Z+D$, one has for every integer $j\ge0$
\begin{equation}\label{eq:induction-ranks}
 \dim_{\C(S)}H^0\!\left(Y_\eta,
   \left(K_{Y/S}+\pi^*(J+jkL)\right)|_{Y_\eta}\right)=1.
\end{equation}
\end{enumerate}
If $S$ is a point, assertion \textup{(i)} already gives a nonzero
section of $kL$ on $Z$.
\end{lemma}

\begin{proof}
Consider the nonempty complete linear systems of
\begin{equation}\label{eq:systems}
              B=uL+vM+wD,
       \qquad u,v,w\in\mathbb Z_{\ge0}.
\end{equation}
Their degrees against $\alpha$ are nonpositive. A divisor whose
complete system defines a generically finite map is big. But a big
class is an interior point of the pseudoeffective cone, and so pairs
strictly positively with every nonzero member of the dual cone.
Thus every system in \eqref{eq:systems} has image dimension less than
$d$.

Choose $B_0$ in this collection with maximal image dimension $s$, and
let $V$ be the closed image of its rational map. This maximum exists
because the possible dimensions are integers between $0$ and $d-1$.
Let $K$ be the relative algebraic closure of $\C(V)$ in $\C(Z)$;
it is a finite extension of $\C(V)$, since $\C(Z)$ is a finitely
generated field extension. Take a smooth projective model $S$ of $K$,
with a morphism to $V$ after resolution. Its dimension is $s<d$, and
the rational map $Z\dashrightarrow S$ has connected general fiber.

Every ratio of two nonzero sections in any one system
\eqref{eq:systems} belongs to $K$. Indeed, combine that system with
$|B_0|$ by taking products of sections. The resulting linear subsystem
of $|B+B_0|$ gives the joint rational map. If its section ratio were
transcendental over $\C(V)$, the joint image would have dimension
at least $s+1$, contrary to the choice of $B_0$. The ratio is therefore
algebraic over $\C(V)$ and lies in its relative algebraic closure.
We will use this \emph{section-ratio property} both for interpolation
and for rank one.

Resolve $|B_0|$ before applying Lemma~\ref{lem:prepared-model}.
The resulting smooth model carries an effective divisor $R_0$ and a
line bundle $H_0$ on $S$ such that
\begin{equation}\label{eq:fixed-moving}
                  \pi^*B_0=f^*H_0+R_0.
\end{equation}
Here $H_0$ is the pullback of $\OO_V(1)$. If $s>0$ it is big,
since $S\to V$ is generically finite. All additional modifications
preserve the displayed effective decomposition by pullback.

Set $k=m_2-m_1>0$ and take
$\rho=\bar\pi^*(s_2/s_1)$. For $i>2$, the numerator and denominator
in
\begin{equation}\label{eq:interpolation}
                  \frac{s_i^k s_1^{m_i-m_2}}
                       {s_2^{m_i-m_1}}
\end{equation}
are sections of the same line bundle of the form
\eqref{eq:systems}: both have $M$ coefficient $m_i-m_1$ and $L$
coefficient $m_2(m_i-m_1)$. Thus \eqref{eq:interpolation} is a
nonzero element of $K$. It has order zero at every horizontal prime
$Q$ of $\bar Y$, so
\begin{equation}\label{eq:horizontal-interpolation}
 k\bigl(\ord_Q(\bar\pi^*s_i)-\ord_Q(\bar\pi^*s_1)\bigr)
        =(m_i-m_1)\ord_Q(\rho).
\end{equation}
The left side is bounded below independently of $i$, since each
$\bar\pi^*s_i$ is regular. As $m_i\to\infty$, the right side could
have this lower bound only if $\ord_Q(\rho)\ge0$. This proves
\textup{(i)}. If $S$ is a point, every prime is horizontal;
normality makes $\rho$ regular, and birational pushdown gives the
claimed section on $Z$. In this case the rank assertion also follows
from the section-ratio property. Indeed, the adjoint space in
\eqref{eq:induction-ranks} identifies with $H^0(Z,D+jkL)$:
the Jacobian divisor is effective and exceptional, so a rational
section allowed to have poles only there extends downstairs across
codimension two on smooth $Z$. The section $s_D\rho^j$ is nonzero,
and any two sections of $D+jkL$ have ratio in $K=\C$ by
\eqref{eq:systems}. Hence this space has dimension one for every
$j\ge0$, including $j=0$.

Assume $s>0$. We now verify \eqref{eq:induction-ranks} in every
nonnegative degree. The identity
\begin{equation}\label{eq:adjoint-identification}
 K_{Y/S}+\pi^*(J+jkL)
       =K_{Y/Z}+\pi^*(D+jkL)-f^*K_S
\end{equation}
shows which exceptional divisor is involved. Over the generic point
of $S$, a base canonical frame is immaterial. The Jacobian section,
$s_D$, and $\rho^j$ give a nonzero section of the right side.
To justify regularity, \textup{(i)} implies that $\rho$ is regular
on $\bar Y$ over a dense base open: remove the images of its finitely
many vertical polar divisors. Its pullback to $Y$ is then regular
there. Thus the adjoint space has dimension at least one, also for
$j=0$.

Take any two nonzero sections of that generic-fiber space. Choose a
rational frame of $K_S$ and regard them, by
\eqref{eq:adjoint-identification}, as rational sections of
$K_{Y/Z}+\pi^*(D+jkL)$ with no horizontal poles. Their finitely
many vertical polar divisors can be cleared by one effective Cartier
divisor $C$ on $S$: enclose their images, including any of codimension
at least two, in a Cartier divisor and increase its multiplicities so
that $f^*C$ dominates both polar divisors. Since $H_0$ is big, there
are an integer $a>0$ and a nonzero section of $aH_0-C$.
Multiplication by this section and by $s_C$ makes both rational
sections regular after the common base twist $af^*H_0$.

Now multiply both by the canonical section of $aR_0$. Formula
\eqref{eq:fixed-moving} produces two regular sections of
\[
                   K_{Y/Z}+\pi^*(D+jkL+aB_0)
\]
with the same ratio as the original pair. They descend to sections
of $D+jkL+aB_0$ on $Z$. In detail, $K_{Y/Z}$ is effective and
exceptional, and $\pi_*\OO_Y(K_{Y/Z})=\OO_Z$: a rational function
with poles only on exceptional divisors is regular away from a
codimension-two subset of smooth $Z$ and hence regular everywhere.
Apply this fact after trivializing the line bundle on $Z$.
The descended line belongs to \eqref{eq:systems}; the section-ratio
property therefore places their ratio in $K=\C(S)$.
Any two nonzero vectors in the generic-fiber section space are thus
proportional over its ground field. This proves dimension one for
all $j\ge0$, without discarding the zeroth degree.
\end{proof}

The same maximal system supplied both inputs for transfer. The
interpolation identity removed horizontal poles, while the big line
$H_0$ on the chosen base let us clear arbitrary vertical poles in
the rank calculation. No equidimensionality of the smooth resolution
was used: poles over higher-codimension base sets were cleared before
pushing down the exceptional Jacobian.

\subsection{Induction on dimension}

\begin{proof}[Proof of Theorem~\ref{thm:finite-volume}]
We induct on $d=\dim Z$. In dimension zero $Z$ is a point, every
line bundle is trivial, and the conclusion holds with $a=1,b=0$.
Assume $d>0$ and suppose that no $aL+bD$, with integers $a>0$ and
$b\ge0$, has a section.

The bundle $L$ is pseudoeffective because it has a semipositive metric.
The class of $L+D$ is therefore pseudoeffective. It is not big:
otherwise some positive multiple of $L+D$ would have a nonzero
section. Corollary~\ref{h:slope:face} gives a nonzero movable class
$\alpha$ with $L\cdot\alpha=D\cdot\alpha=0$ and
$\mu_{\min,\alpha}(T_Z)\ge0$. Apply Lemma~\ref{lem:determinant}
to obtain the fixed line $M$ and unbounded sections. Then apply
Lemma~\ref{lem:maximal-base}.

If its base is a point, the resulting section contradicts our
assumption. Otherwise Proposition~\ref{prop:transfer}, applied to
that prepared diagram and the two original metrics, supplies a line
$A$ and an effective integral divisor $D_S$ on the smaller smooth
projective base $S$. The line $A$ has a semipositive metric with
locally bounded weights, while $-K_S+D_S$ has a semipositive metric
with finite distinguished adjoint volume.

The no-forms condition also passes to $S$. A holomorphic form on
$S$ pulls back along the dominant rational map on its domain of
definition. Properness of $S$ extends the map at the generic point
of every divisor of $Z$, so that domain can be taken to have
codimension-two complement. The pulled-back form extends across that
complement because $\Omega_Z^q$ is locally free. Dominance in
characteristic zero makes pullback injective. Hence
$H^0(S,\Omega_S^q)=0$ for every $q>0$.

Since $\dim S<d$, induction produces integers $a>0,b\ge0$ and a
nonzero section of $aA+bD_S$. The lifting conclusion of
Proposition~\ref{prop:transfer} gives a nonzero section of
$akL+cD$ on $Z$ for some $c\ge0$. Its exponent of $L$ is positive,
so this contradicts the assumption and completes the induction.
\end{proof}
